\CharacterSheet{buttler}{Mordoma}{

\shortdescription{A administradora da casa, antiga comerciante de olhar atento.}

\ObjectivesHeading{Objetivos}
\begin{itemize}
  \item Receber dinheiro em trocas comerciais. Por exemplo, vendendo informações. (\PositiveReward{1}, \NegativeReward{0})
\end{itemize}

\LosingConditionsHeading{Condições de derrota}
Não conseguir que quem matou \Character{newcomer} seja expulso por votação, caso haja uma pessoa responsável pelo homicídio.

\SecretsHeading{Segredos}
\begin{itemize}
  \item Não posso deixar que descubram as ocasiões em que recebi dinheiro em troca de serviços. (\NegativeReward{1}, \PositiveReward{1}) por cada uma dessas ocasiões
\end{itemize}

\ActionsHeading{Ações}
\begin{itemize}
  \item Mantém a casa a funcionar. Oferece ajuda, como faria uma boa mordoma.
\end{itemize}

\BackgroundHeading{Contexto}

\timestamp{Décadas depois da morte de Cristo}, era uma comerciante conhecida por manter as contas em ordem e fazer com que os negócios difíceis parecessem não dar trabalho. \Character{vampirehead} convidou-me para a sua família porque precisava dos meus talentos de gestão. Aceitei porque a imortalidade é o melhor investimento que já fiz. Deram-me uma festa só com \Character{vampirehead}, \Character{enforcer} e \Character{usurper}. Foi horrível; precisavam mesmo de mim.

Desde então, mantenho este castelo abastecido, organizado e financeiramente estável.

\timestamp{No ano 100 d.C.}, foi acolhido um novo membro, \Character{ghoul}, uma espécie de artista. Foi apresentado por \Character{enforcer}, que parecia querer uma relação mais profunda com ele. Na festa, reparei que a ``poção proibida'' oferecida por \Character{vampirehead} parecia ter uma tonalidade diferente. Quando lhe perguntei, ofereceu-me dinheiro pelo meu silêncio, pelo que não insisti no assunto.

Ao longo dos séculos seguintes, a família cresceu. Cada chegada trazia outro quarto para preparar, novos hábitos a acomodar e outra fonte de atrito que alguém tinha de resolver antes do jantar.

\timestamp{No ano 500 d.C.}, chegou \Character{supremacist}. Tinha sido governante e esperava que a casa o reconhecesse. Era uma companhia exigente, mas compreendia o cerimonial e tratava a família com mais respeito do que a maioria dos humanos alguma vez tinha mostrado. Quando \Character{vampirehead} lhe atribuiu a tarefa de avaliar possíveis novos membros, aprendi a guardar para mim as opiniões sobre os potenciais convidados.

\timestamp{No ano 700 d.C.}, \Character{lover} juntou-se a nós como protetor da família. Era um guerreiro errante, feito para soluções diretas e pouco interessado na organização doméstica que as sustentava. Ainda assim, era cortês comigo, e é útil ter alguém capaz de pôr fim a uma discussão antes que saia cara.

\timestamp{No ano 1000 d.C.}, \Character{prevlover} tornou-se o nosso investigador. Tinha o hábito de reparar em pormenores que os outros esperavam manter privados, o que fazia dele um convidado inconveniente, mas um familiar útil.

\timestamp{No ano 1400 d.C.}, \Character{traitor} foi admitida depois de \Character{lover} ter insistido. Era conhecida por ajudar as pessoas, o que a tornou, desde o início, uma escolha invulgar para esta família.

\timestamp{Há dois anos}, \Character{vampirehead} apresentou um humano à família numa festa: \Character{newcomer}. Foi uma boa ajuda nas tarefas depois da festa; uma mudança agradável.

\timestamp{Há dois meses}, como era costume, \Character{vampirehead} marcou \timestamp{16/7/2026} como data da festa para acolher \Character{newcomer} na família. Todos os vampiros tinham de vir.

\timestamp{Há umas duas semanas}, \Character{vampirehead} perguntou se mais alguém, além dela, de mim e de \Character{newcomer}, queria ajudar. Disse que \Character{usurper} se tinha oferecido e chegaria ao castelo duas noites antes da festa. A lista final ficou com \Character{vampirehead}, \Character{newcomer}, \Character{usurper} e, como é habitual, comigo.

\timestamp{Duas noites antes da festa}, \Character{usurper} e \Character{newcomer} chegaram ao castelo para ajudar nos preparativos. \Character{vampirehead} e eu já lá estávamos.

Não consigo reconstituir toda a \timestamp{manhã de ontem}. Pelos meus registos atuais, suspeito de que aceitei um pagamento e lembro-me de que tinha sido posta de parte uma garrafa para \Character{newcomer}.

\timestamp{Ontem, às 12:00}, almocei com \Character{vampirehead}, \Character{newcomer} e \Character{usurper}. Estava prestes a servir o vinho em questão quando \Character{usurper} interveio. Disse que tinha trazido uma garrafa melhor para cada um de nós, e tinha razão. \Character{usurper} vem sempre em nosso socorro. Lembro-me de que queria servir a garrafa que tinha, mas não quis contrariar \Character{usurper}; podia sempre oferecê-la mais tarde.

\timestamp{Desde as 16:00 de ontem}, tenho notado problemas de memória. Não é só ter-me esquecido de coisas; também me custa guardar novas recordações.

\timestamp{Hoje de manhã, às 11:00}, vi alguém a falar com \Character{newcomer} em privado lá fora. Sei que aceitei dinheiro pelo meu silêncio, mas não me lembro de mais pormenores.

\GMNotes{
\paragraph{Porque se pode suspeitar desta personagem?}
\Character{buttler} tratou do vinho e aceitou dinheiro para guardar segredos. É fácil imaginar alguém a pagar-lhe para ajudar no homicídio.

\paragraph{Como se pode provar a inocência desta personagem?}
\Character{vampirehead} também viu \Character{usurper} substituir o vinho ao almoço. As notas sobre a poção explicam a importância dessa garrafa.

\paragraph{Qual é o papel desta personagem na história e no jogo?}
Tem informações úteis e quer que lhe paguem por elas. As falhas de memória tornam alguns pormenores mais difíceis de recuperar.

\paragraph{Como pode esta personagem encaminhar a investigação para a solução?}
Lembra-se de \Character{usurper} substituir o vinho reservado para \Character{newcomer}. É uma das ligações mais claras entre \Character{usurper} e a primeira dose.
}
}
